% Makeup final examination, question-only version.
% Compile make_up_exam_answers.tex to produce the version with solutions.

\ifdefined\MakeupAnswerVersion
  \documentclass[12pt,a4paper,answers,addpoints]{exam}
\else
  \documentclass[12pt,a4paper,addpoints]{exam}
\fi

\def\CourseTitle{Computer Programming}
\def\ExamName{Final Examination}
\def\CourseNumber{BE104}
\def\ExamDate{January, 5th, 2027}
\def\Duration{2 Hours}
\def\ExamTime{15:00--17:00}
\def\Term{2609}
\def\ExamDuration{2 Hours}

\usepackage[left=2cm, right=2cm, top=2cm, bottom=2cm]{geometry}

\usepackage{amsmath}
\usepackage{amsthm}
\usepackage{amsfonts}
\usepackage{amssymb}
\usepackage{mathrsfs}
\usepackage{tasks}
% \usepackage[margin=1in]{geometry}
\checkboxchar{$\Box$}
\usepackage{setspace}
%\usepackage{bbding}
%\usepackage{pifont}
\usepackage{wasysym}
\usepackage[shortlabels]{enumitem}
\renewcommand{\qedsymbol}{$\blacksquare$}
\usepackage{array} 
\usepackage{multicol}
\setlength{\columnsep}{1cm}
\usepackage{afterpage}
\usepackage{color} 
\usepackage{graphicx} % in preamble
\usepackage{listings}
\usepackage{upquote}

\lstset{
    escapeinside={(*@}{@*)},           % Allow LaTeX commands between (*@ and @*)
}

% new macros
\newcommand{\todo}[1]{\textcolor{red}{(To do): #1}}
\newcommand{\tored}[1]{\textcolor{red}{\textbf{#1}}}
\usepackage{xeCJK}
\defaultfontfeatures{Ligatures=TeX}

\setCJKmainfont[
    BoldFont={Noto Sans TC}
]{Noto Sans TC}

% Required by listings/typewriter text
\setCJKmonofont{Noto Sans TC}

\newcommand\blankpage{%
    \null
    \thispagestyle{empty}%
    \addtocounter{page}{-1}%
    \shownoanswers{\newpage}}
%% https://tex.stackexchange.com/a/203986
% Create a True-False question format
\newcommand*{\TrueFalse}[1]{%
\ifprintanswers
    \ifthenelse{\equal{#1}{T}}{%
        \textbf{TRUE}\hspace*{14pt}False
    }{
        True\hspace*{14pt}\textbf{FALSE}
    }
\else
    {True}\hspace*{20pt}False
\fi
} 
%% The following code is based on an answer by Gonzalo Medina
%% https://tex.stackexchange.com/a/13106/39194
\newlength\TFlengthA
\newlength\TFlengthB
\settowidth\TFlengthA{\hspace*{1.16in}}
\newcommand\TFQuestion[2]{%
    \setlength\TFlengthB{\linewidth}
    \addtolength\TFlengthB{-\TFlengthA}
    \parbox[t]{\TFlengthA}{\TrueFalse{#1}}\parbox[t]{\TFlengthB}{#2}}
%%

%%
\pagestyle{headandfoot}
%\extraheadheight{.25in}
\runningheader{{\CourseNumber} - {\CourseTitle}}
{}
{{\ExamName} - {\Term}} 
\firstpagefooter{}{Page \thepage\ of \numpages}{Please go on to the next page\ldots} %
\runningfooter{}{Page \thepage\ of \numpages}{\iflastpage{End of exam.}{Please go on to the next page\ldots}}

\renewcommand{\fillinlinelength}{1cm}  % change default length

\renewenvironment{solution}
  {\par\noindent\textbf{Solution: }\color{red}} % begin code
  {\par}    

% --- show solution chunks only when 'answers' is active ---
\usepackage{xcolor} % for \color

% Inline macro: use as \showsol{<solution text>}
\newcommand{\showsol}[1]{\ifprintanswers{\textbf{\color{red}\underline{Solution:} }\color{red}#1}\fi}

\newenvironment{show_code_sol}
  {\par\noindent\textbf{Solution:}\par
   \lstset{basicstyle=\ttfamily\small\color{red}}}
  {\par}


\usepackage{listings}
\usepackage{comment} % for \excludecomment (safe with verbatim)
% \usepackage[T1]{fontenc}


% Environment version that mimics 'solution' but guarded
\newenvironment{solonly}
  {\ifprintanswers\par\noindent\textbf{Solution: }\color{red}}
  {\par\fi}

% If you want to keep using the built-in 'solution' name,
% make sure it stays conditional:
\renewenvironment{solution}
  {\ifprintanswers\par\noindent\textbf{Solution: }\color{red}}
  {\par\fi}

\newcommand{\shownoanswers}[1]{\ifprintanswers\else #1\fi}

% Base style for C
% One listings + comment load
\usepackage{listings}
\usepackage{comment} % safe with verbatim

% --- styles (no frame, no line numbers) ---
\lstdefinestyle{examc}{
  language=C,
  basicstyle=\ttfamily\small,
  numbers=none,
  frame=none,
  columns=fullflexible,
  keepspaces=true,
  showstringspaces=false,
  tabsize=4,
  breaklines=true,
  keywordstyle=\bfseries,
  commentstyle=\itshape,
  stringstyle=\ttfamily
}
\lstdefinestyle{examc-sol}{
  style=examc,
  basicstyle=\ttfamily\small\color{red}
}

% --- Environments ---
\lstnewenvironment{Ccode}[1][]
  {\lstset{style=examc,#1}}
  {}

% Define CcodeSol only when answers are printed; otherwise exclude it.
\ifprintanswers
  \lstnewenvironment{CcodeSol}[1][]
    {\lstset{style=examc-sol,#1}}
    {}
\else
  \excludecomment{CcodeSol} % body is ignored; no verbatim processing
\fi

% Inline C
\newcommand{\Cinline}[1]{\lstinline[language=C]{#1}}

% ---------------------------------------------------------------------
  


\begin{document}

% \noindent
% \begin{minipage}[t]{0.48\textwidth}
%     \vspace{-10pt}
%     \fbox{\parbox[c][1.2cm][c]{4cm}{Seat\\No.}}
% \end{minipage}

% \vspace{-2cm}
% \begin{center}
%     \includegraphics[width=0.27\textwidth]{must-logo.pdf}
% \end{center}

% \begin{center}
%     \vspace*{0.5cm}
%     {\textbf{\Huge 澳\hspace{2em}門\hspace{2em}科\hspace{2em}技\hspace{2em}大\hspace{2em}學}}\\[0.3cm]
%     {\large \textbf{MACAU UNIVERSITY OF SCIENCE AND TECHNOLOGY}}\\[1cm]

%     {\Large \textbf{School of Computer Science and Engineering}}\\[0.2cm]
%     {\Large \textbf{Faculty of Innovation Engineering}}\\[0.2cm]
%     {\Large \textbf{Bachelor of Science}}\\[0.5cm]

%     {\large \textbf{2609}}\\[0.6cm]
%     {\Large \textbf{\CourseNumber \quad \CourseTitle}}\\[0.6cm]
% \end{center}

% \begin{center}
% \begin{tabular}{@{}p{4cm}p{8cm}@{}}
%     \textbf{No. of Pages:} & {\numpages} Pages (including this page)\\[0.35cm]
%     \textbf{Exam Date:} & \ExamDate\\[0.35cm]
%     \textbf{Exam Duration:} & \ExamDuration\ (\ExamTime)\\[0.35cm]
%     \textbf{Student Name:} & \underline{\hspace{7cm}}\\[0.35cm]
%     \textbf{Student Number:} & \underline{\hspace{7cm}}\\
% \end{tabular}
% \end{center}

% \vspace{0.2cm}

% \noindent
% \textbf{Please read the following instructions carefully.}
% \begin{itemize}[leftmargin=*,itemsep=2pt,topsep=3pt]
%     \item This question booklet contains \textbf{\numquestions} questions and
%     \textbf{\numpages} pages (including this page), for a total of
%     \textbf{\numpoints} points/marks. Check that all pages are present.
%     \item Answer all questions in the spaces provided. For each multiple-choice
%     question, select only one answer.
%     \item Return this Question Booklet, together with all extra unused answer
%     sheets and draft paper, at the end of the examination.
%     \item This is a \textbf{Closed Book Exam}. Dictionaries, calculators, and
%     electronic devices are \textbf{NOT permitted}.
%     \item \textbf{Answers must be provided in English to receive marks.}
% \end{itemize}

% \vfill

% \hspace{0.5in}
\begin{questions}
\addpoints

\newpage
\fullwidth{\section*{Section I: Multiple-Choice Questions (20 points)}}
% \fullwidth{Select the one best answer for each question. Each question is worth 2 points.}

\question[2] Given \texttt{int value = 25;}, which statement correctly declares
the pointer \texttt{p} and initializes it with the address of \texttt{value}?
\begin{choices}
    \choice \texttt{int p = \&value;}
    \choice \texttt{int *p = value;}
    \choice \texttt{int *p = \&value;}
    \choice \texttt{int \&p = value;}
\end{choices}
\showcorrectanswer{C. \texttt{int *p = \&value;}}
\showoutcome{1}{b (b2)}

\question[2] The pointer \texttt{q} points to the first element of an integer array
\texttt{data}. Which expression represents \texttt{data[3]}?
\begin{oneparchoices}
    \choice \texttt{*q + 3}
    \choice \texttt{\&q + 3}
    \choice \texttt{*(q + 3)}
    \choice \texttt{*q[3]}
\end{oneparchoices}
\showcorrectanswer{C. \texttt{*(q + 3)}}
\showoutcome{1}{b (b2)}

\question[2] Given the declarations below, which expression accesses the score
through the pointer \texttt{sp}?
\begin{lstlisting}[style=examc]
struct Student {
    char name[20];
    int score;
};
struct Student s;
struct Student *sp = &s;
\end{lstlisting}
\begin{oneparchoices}
    \choice \texttt{sp.score}
    \choice \texttt{sp->score}
    \choice \texttt{*sp.score}
    \choice \texttt{score.sp}
\end{oneparchoices}
\showcorrectanswer{B. \texttt{sp->score}}
\showoutcome{1}{b (b2)}

\question[2] Which statement correctly reads a \texttt{double} from the file
represented by \texttt{fp} and stores it in \texttt{value}?
\begin{choices}
    \choice \texttt{fscanf("\%lf", fp, value);}
    \choice \texttt{fscanf(fp, "\%lf", value);}
    \choice \texttt{fscanf(fp, "\%lf", \&value);}
    \choice \texttt{fscanf(\&fp, "\%lf", \&value);}
\end{choices}
\showcorrectanswer{C. \texttt{fscanf(fp, "\%lf", \&value);}}
\showoutcome{1}{b (b2)}

\question[2] What is printed by the following program fragment?
\begin{lstlisting}[style=examc]
int total = 0;
for (int i = 1; i <= 3; i++)
    total += i * 2;
printf("%d", total);
\end{lstlisting}
\begin{oneparchoices}
    \choice \texttt{6}
    \choice \texttt{9}
    \choice \texttt{12}
    \choice \texttt{18}
\end{oneparchoices}
\showcorrectanswer{C. \texttt{12}}
\showoutcome{2}{b (b2)}

\question[2] Which statement correctly reads an integer into \texttt{count} and a
character into \texttt{grade}?
\begin{choices}
    \choice \texttt{scanf("\%d \%c", count, grade);}
    \choice \texttt{scanf("\%d \%c", \&count, \&grade);}
    \choice \texttt{scanf(\%d \%c, \&count, \&grade);}
    \choice \texttt{scanf("\%d \%c", *count, *grade);}
\end{choices}
\showcorrectanswer{B. \texttt{scanf("\%d \%c", \&count, \&grade);}}
\showoutcome{1}{b (b2)}

\question[2] Which statement about \texttt{break} and \texttt{continue} inside a
loop is correct?
\begin{choices}
    \choice \texttt{break} skips only the current iteration.
    \choice \texttt{continue} immediately terminates the loop.
    \choice \texttt{break} exits the loop, while \texttt{continue} proceeds
    to the next iteration.
    \choice Both statements always terminate the program.
\end{choices}
\showcorrectanswer{C. \texttt{break} exits the loop, while \texttt{continue}
proceeds to the next iteration.}
\showoutcome{1}{b (b2)}

\question[2] What is printed by the following program fragment?
\begin{lstlisting}[style=examc]
int code = 1;
switch (code) {
    case 0: printf("A"); break;
    case 1: printf("X");
    case 2: printf("Y"); break;
    default: printf("Z");
}
\end{lstlisting}
\begin{oneparchoices}
    \choice \texttt{X}
    \choice \texttt{XY}
    \choice \texttt{XYZ}
    \choice \texttt{Z}
\end{oneparchoices}
\showcorrectanswer{B. \texttt{XY}}
\showoutcome{2}{b (b2)}

\question[2] Suppose \texttt{x = 4} and \texttt{y = 9}. Which condition is true?
\begin{choices}
    \choice \texttt{x > y \&\& y > 0}
    \choice \texttt{x < y \&\& y - x == 5}
    \choice \texttt{x == y || x > 5}
    \choice \texttt{x + y < 10}
\end{choices}
\showcorrectanswer{B. \texttt{x < y \&\& y - x == 5}}
\showoutcome{1}{b (b2)}

\question[2] Which of the following is a valid C variable name?
\begin{choices}
    \choice \texttt{3level}
    \choice \texttt{student-name}
    \choice \texttt{\_level3}
    \choice \texttt{while}
\end{choices}
\showcorrectanswer{C. \texttt{\_level3}}
\showoutcome{1}{b (b2)}

\newpage
\fullwidth{\section*{Section II: Write Program Output (30 points)}}

\question[10] Write the exact output of the following program, including spaces.
Then briefly identify which array element each printed expression represents.
\begin{lstlisting}[style=examc]
#include <stdio.h>

int main(void) {
    int a[8] = {3, 6, 9, 12, 15, 18, 21, 24};
    int *p = a;
    int *ap[2] = {&a[6], &a[2]};

    printf("%d %d %d %d\n", a[7], *(p + 4), *ap[0], *ap[1]);
    return 0;
}
\end{lstlisting}
\shownoanswers{\vspace{4cm}}
\showsol{The exact output is \texttt{24 15 21 9}. The four expressions represent
\texttt{a[7]}, \texttt{a[4]}, \texttt{a[6]}, and \texttt{a[2]}, respectively.

\textbf{Marking:} 2 points for each correct value and 2 points for a correct
identification or explanation of the expressions.}
\showoutcome{2}{b (b2)}

\newpage
\question Consider the following two functions.
\begin{lstlisting}[style=examc]
char element_index(const char a[], int n) {
    return a[n];
}

char element_pointer(const char *p, int n) {
    return *(p + n);
}
\end{lstlisting}
Let the character array be:
\begin{lstlisting}[style=examc]
char letters[] = {'C', 'O', 'D', 'I', 'N', 'G', '\0'};
\end{lstlisting}
\begin{parts}
    \part[3] What is printed by the following statement?
\begin{lstlisting}[style=examc]
printf("%c", element_index(letters, 3));
\end{lstlisting}
    \shownoanswers{\vspace{4cm}}
    \showsol{\texttt{I}}
    \showoutcome{2}{b (b2)}

    \part[3] What is printed by the following statement?
\begin{lstlisting}[style=examc]
printf("%c", element_pointer(letters, 5));
\end{lstlisting}
    \shownoanswers{\vspace{4cm}}
    \showsol{\texttt{G}}
    \showoutcome{2}{b (b2)}

    \part[4] In your words, what do the two functions do?
    \shownoanswers{\vspace{4cm}}
    \showsol{Both functions return the element at index \texttt{n} of the supplied
    character array. The first uses array indexing; the second uses pointer arithmetic.}
    \showoutcome{1}{b (b2)}
\end{parts}
\par\medskip

\newpage
\question Study the following function. Positions are counted starting from 1.
\par\vspace{0.5em}
\begin{lstlisting}[style=examc]
void print_odd_positions(const char s[]) {
    int i = 0;
    while (s[i] != '\0') {
        putchar(s[i]);
        i += 2;
    }
    putchar('\n');
}
\end{lstlisting}
\begin{parts}
    \part[3] What is printed by \texttt{print\_odd\_positions("ABCDEFG");}?
    \shownoanswers{\vspace{1.5cm}}
    \showsol{\texttt{ACEG}}
    \showoutcome{2}{b (b2)}

    \part[3] What is printed by \texttt{print\_odd\_positions("Computer");}?
    \shownoanswers{\vspace{1.5cm}}
    \showsol{\texttt{Cmue}}
    \showoutcome{2}{b (b2)}

    \part[4] Explain what \texttt{print\_odd\_positions()} does.
    \shownoanswers{\vspace{3cm}}
    \showsol{The function prints the characters in the 1st, 3rd, 5th, and subsequent
    odd-numbered positions, followed by a newline.}
    \showoutcome{2}{b (b2)}
\end{parts}

\newpage
\fullwidth{\section*{Section III: Fill in the Blanks (30 points)}}

\question[10] Complete the five blanks so that the program uses a \texttt{while}
loop to count and print the number of lowercase English letters in a string. 
\begin{lstlisting}[style=examc]
#include <stdio.h>

int count_lowercase((*@\fillin[\tored{const char *p}][3cm]@*)) {
    int count = 0;
    while ((*@\fillin[\tored{*p != '\textbackslash0'}][3cm]@*)) {
        if ((*@\fillin[\tored{*p >= 'a' \&\& *p <= 'z'}][4.5cm]@*))
            count++;
        (*@\fillin[\tored{p++}][2cm]@*);
    }
    return count;
}

int main(void) {
    char text[] = "C Programming 2027";
    printf("Lowercase letters = %d\n",
           (*@\fillin[\tored{count\_lowercase(text)}][4cm]@*));
    return 0;
}
\end{lstlisting}
\showsol{The completed program prints \texttt{Lowercase letters = 10}.}
\showoutcome{2}{b (b2)}

\question[10] Complete the three blanks so that the function sorts the array in
descending order using bubble sort. 
\begin{lstlisting}[style=examc]
void bubble_sort_descending(int a[], int n) {
    for (int i = 0; (*@\fillin[\tored{i < n - 1}][3.2cm]@*); i++) {
        for (int j = 0; (*@\fillin[\tored{j < n - 1 - i}][3.2cm]@*); j++) {
            if ((*@\fillin[\tored{a[j] < a[j + 1]}][3.8cm]@*)) {
                int temp = a[j];
                a[j] = a[j + 1];
                a[j + 1] = temp;
            }
        }
    }
}
\end{lstlisting}
\showsol{The loops perform bubble-sort passes, and the comparison swaps adjacent
elements when the left value is smaller, producing descending order.}
\showoutcome{2}{b (b2)}

\newpage
\question[10] Complete the three blanks so that the function \textbf{uses pointer
arithmetic} to print the integers in the 2nd, 4th, 6th, and subsequent even-numbered
positions of the array. Positions are counted starting from 1. 
{\color{red}Using array indexing instead of pointer arithmetic may not receive full score.}
\begin{lstlisting}[style=examc]
#include <stdio.h>

void print_even_positions(const int *p, int n) {
    for (int i = (*@\fillin[\tored{1}][2cm]@*); i < n;
         (*@\fillin[\tored{i += 2}][2.5cm]@*)) {
        printf("%d ", (*@\fillin[\tored{*(p + i)}][3cm]@*));
    }
    putchar('\n');
}

int main(void) {
    int a[8] = {5, 10, 15, 20, 25, 30, 35, 40};
    print_even_positions(a, 8);
    return 0;
}
\end{lstlisting}
\showsol{The program prints \texttt{10 20 30 40}.}
\showoutcome{2}{b (b2)}

\newpage
\fullwidth{\section*{Section IV: Write C Code (20 points)}}

\question[10] The following program stores cars of different brands and models in
an array of structures. Complete the function \texttt{total\_price()} so that it
returns the sum of all car prices. The program must print the total price.
\begin{lstlisting}[style=examc]
#include <stdio.h>

struct Car {
    char brand[30];
    char model[30];
    double price;
};

double total_price(const struct Car *p, int n) {
    /* Write your code here. */
}

int main(void) {
    struct Car cars[] = {
        {"Toyota", "Corolla", 23000.0},
        {"Honda", "Civic", 25000.0},
        {"Tesla", "Model 3", 40000.0},
        {"BMW", "X1", 42000.0}
    };
    int n = sizeof(cars) / sizeof(cars[0]);

    printf("Total price: %.2f\n", total_price(cars, n));
    return 0;
}
\end{lstlisting}
\shownoanswers{\vspace{8cm}}
\begin{CcodeSol}
double total_price(const struct Car *p, int n) {
    double total = 0.0;
    for (int i = 0; i < n; i++)
        total += p[i].price;
    return total;
}
\end{CcodeSol}
\showsol{The output is \texttt{Total price: 130000.00}.

\textbf{Marking:} initialization of the total (2 points), correct loop (2 points),
correct accumulation of each price (4 points), and correct return value (2 points).}
\showoutcome{2}{b (b2)}

\newpage
\question[10] Complete both functions below so that each returns the sum of the
integers from 1 to \texttt{n}. The function \texttt{sum\_loop()} must use a loop,
and \texttt{sum\_recursive()} must use recursion. 
\begin{lstlisting}[style=examc]
#include <stdio.h>

int sum_loop(int n) {
    /* Write your code here. */
}

int sum_recursive(int n) {
    /* Write your code here. */
}

int main(void) {
    int n = 10;
    printf("Loop sum = %d\n", sum_loop(n));
    printf("Recursive sum = %d\n", sum_recursive(n));
    return 0;
}
\end{lstlisting}
\shownoanswers{\vspace{8cm}}
\begin{CcodeSol}
int sum_loop(int n) {
    int total = 0;
    for (int i = 1; i <= n; i++)
        total += i;
    return total;
}

int sum_recursive(int n) {
    if (n <= 0)
        return 0;
    return n + sum_recursive(n - 1);
}
\end{CcodeSol}
\showsol{The program prints \texttt{Loop sum = 55} and
\texttt{Recursive sum = 55}. Equivalent loop and recursive solutions receive full credit.}
\showoutcome{2}{b (b2)}

\end{questions}

\vfill
\noindent\makebox[2.5in]{\textbf{\Large Score:} \hrulefill}

\end{document}
